\section{Labelled determinants and the numerical height budget}
\label{sec:sr-labelled-height-budget}

This chapter proves the determinant and numerical deductions in the
height comparison.  The hypotheses about an actual integral model,
its labelled lattices and its Hodge metric are stated explicitly.
In particular, the deductions below do not construct that model,
prove its Hodge--Tate freeness, or establish its flatness comparison.

Write \(d_E=[E:\mathbb Q]\), \(d_L=[L:\mathbb Q]\), and
\(R_E=|\operatorname{disc}E|^{1/d_E}\).
All finite indices are indices of additive groups.
For two full lattices \(M,M'\) in the same vector space put
\[
 \iota(M,M')=
 \log\#\bigl(M/(M\cap M')\bigr)
 -\log\#\bigl(M'/(M\cap M')\bigr).
 \tag{HB1}\label{eq:sr-hb-signed-index}
\]
Thus \(\iota(M,M')\) is negative when \(M\) is a proper
sublattice of \(M'\).  The sign will matter in the torus cancellation.

\subsection{Determinant identities over integer rings}

Let \(M\) be a full \(\mathcal O_L\)-lattice of rank \(n\)
in an \(L\)-space, equipped with conjugation-compatible positive
Hermitian forms at the complex embeddings.
For an \(L\)-basis \(b=(b_1,\ldots,b_n)\), put
\(M_b=\sum_i\mathcal O_Lb_i\).  Then
\[
 \widehat{\deg}_n M
 =\frac{\iota(M,M_b)
       -\frac12\sum_{\sigma:L\hookrightarrow\mathbb C}
                 \log\det\operatorname{Gram}_{\sigma}(b)}
        {d_L}.
 \tag{HB2}\label{eq:sr-hb-degree-index}
\]
To prove this without assuming a free global lattice, localize
at a finite prime \(\mathfrak p\).  Both lattices are free over
the DVR \(\mathcal O_{L,\mathfrak p}\).
Smith diagonalization writes the first relative to the second
with elementary divisors \(\varpi^{a_1},\ldots,\varpi^{a_n}\).
Its contribution to \(\iota\) is
\(-\sum_i a_i\log N\mathfrak p\).
These are exactly the valuations of its determinant fractional
ideal, with the sign in the arithmetic degree.
Only finitely many primes contribute.
At infinity the norm of \(b_1\wedge\cdots\wedge b_n\)
is the square root of the Gram determinant.
This proves (HB2), including nonprincipal determinant ideals.
A change of \(b\) cancels between the determinant ideal and
the archimedean term by the product formula.

Suppose a nonzero integral ideal \(\mathfrak a\) satisfies
\(\mathfrak a M\subset F\) and \(\mathfrak a F\subset M\).
At \(\mathfrak p\), the elementary-divisor exponents of \(F\)
relative to \(M\) are between
\(-v_{\mathfrak p}(\mathfrak a)\) and
\(v_{\mathfrak p}(\mathfrak a)\).
Adding their valuations proves
\[
 |\iota(F,M_b)-\iota(M,M_b)|
                    \le n\log N_{L/\mathbb Q}\mathfrak a.
 \tag{HB3}\label{eq:sr-hb-commensurable}
\]
This argument also proves that \(F\) is full and finitely generated:
it lies between \(\mathfrak a M\) and \(\mathfrak a^{-1}M\).
The difference of two signed indices is independent of their
common comparison lattice, as is also immediate prime by prime.

\subsection{The torus determinant, including the different}

Let \(I\) be a torsion-free finite \(\mathcal O_E\)-module of
rank \(r\), embedded in an \(L\)-space \(V\) in such a way that
the induced map \(I\otimes_{\mathbb Z}L\to V\) is injective.
Assume \(L\) contains every embedding of \(E\), and let \(e_s\)
be its rational spectral projectors.  Choose
\(v_1,\ldots,v_r\in I\) such that
\[
 I_0=\sum_i\mathcal O_Ev_i\subset I,\qquad H=[I:I_0]<\infty.
\]
Let \(T=\mathcal O_L I\), and
\(T_s=T\cap e_sV\).  The vectors \(e_sv_i\) form an \(L\)-basis
of \(L T_s=e_s(LI)\); this need not be all of \(e_sV\).  The exact formula is
\[
 \frac{\iota(T_s,\sum_i\mathcal O_Le_sv_i)}{d_L}
                       =-r\log R_E+\frac{\log H}{d_E}.
 \tag{HB4}\label{eq:sr-hb-torus-index}
\]

Here is the integral computation.
For an ideal \(\mathfrak b\) in \(E\), its trace-dual is
\(\mathfrak b^{-1}\mathfrak D_E^{-1}\).
If \(a_j\) is an integral basis of \(\mathfrak b\)
and \(b_j\) its trace-dual basis, then in
\(E\otimes_{\mathbb Q}L\)
\[
                         e_s=\sum_j a_j\otimes s(b_j).
\]
Consequently \(x e_s\) belongs to
\(\mathfrak b\otimes_{\mathbb Z}\mathcal O_L\) precisely when
\(x\,s(\mathfrak b^{-1}\mathfrak D_E^{-1})
\subset\mathcal O_L\), or equivalently
\[
                         x\in s(\mathfrak D_E\mathfrak b)
                                               \mathcal O_L.
 \tag{HB5}
\]
This computation uses the rational projector; it makes no claim
that \(e_s\) is integral at ramified primes.

Choose an ideal pseudo-basis \(I=\bigoplus_i\mathfrak b_i v_i'\).
Apply (HB5) to each summand, and take its determinant.
For every fractional ideal \(\mathfrak c\) of \(E\),
\[
 \log N_{L/\mathbb Q}(s(\mathfrak c)\mathcal O_L)
                    =\frac{d_L}{d_E}\log N_{E/\mathbb Q}\mathfrak c.
\]
Changing from \(v'\) to \(v\) multiplies the determinant
by \(s(\det C)\); its norm is accounted for by the same formula
for principal fractional ideals.
The determinant ideal of \(I\), relative to \(v\), has norm
\(H^{-1}\).  The different has norm
\(|\operatorname{disc}E|\).
The signed-index formula therefore gives
\[
 \iota(T_s,\mathcal O_L e_sv)
 =-\frac{d_L}{d_E}\bigl(r\log|\operatorname{disc}E|-\log H\bigr),
\]
which is (HB4).

\subsection{The labelled degree and its cancellation}

We use the following precise linear data.
There is a full \(\mathcal O_L\)-lattice \(\Lambda\subset V\),
stable under \(\mathcal O_E\), and an \(\mathcal O_E\)-stable
\(\mathbb Z\)-lattice \(I\subset\Lambda\) as above.
There is an \(L\)-linear map \(f:V\to B_L\), with kernel
\(L I\), taking \(\Lambda\) into a Hermitian lattice \(B\).
Its image spans \(B_L\).  The images of distinct \(E\)-eigenspaces
are orthogonal at every infinite embedding.
Put
\[
 \Lambda_s=\Lambda\cap e_sV,\quad
 T_s=(\mathcal O_L I)\cap e_sV,\quad
 M_s=f(\Lambda_s).
\]
Assume the integral kernel is exactly
\(\Lambda\cap LI=\mathcal O_L I\).
Then
\(0\to T_s\to\Lambda_s\to M_s\to0\) is exact.

Choose \(q_1,\ldots,q_k\in\Lambda_s\) such that
\((e_sv_1,\ldots,e_sv_r,q_1,\ldots,q_k)\) is an \(L\)-basis
of \(e_sV\).  Define
\[
 \operatorname{ldeg}_s(F)=\frac1{d_L}
 \left\{\iota\bigl(F,\mathcal O_L(e_sv,q)\bigr)
  -\frac12\sum_\sigma
           \log\det\operatorname{Gram}_{\sigma}(f(q))\right\}.
 \tag{HB6}
\]
At a finite prime choose a basis of the kernel and lift a basis
of the image.  The relative basis matrix is block triangular.
Its determinant is the product of its kernel and quotient
determinants, independently of the chosen lifts.
Summing these DVR identities and applying (HB2) gives
\[
 \operatorname{ldeg}_s(\Lambda_s)
 =\frac{\iota(T_s,\mathcal O_Le_sv)}{d_L}
                              +\widehat{\deg}_n M_s.
 \tag{HB7}
\]
Only the quotient has a Hermitian Gram term in this formula.

Let the chosen labels be indexed by a set of size
\(f_0=d_E/2\).  They need not be all the embeddings.
For each chosen label suppose an ideal \(\mathfrak a_s\ne0\)
satisfies
\[
 \mathfrak a_sF_s\subset\Lambda_s,\quad
 \mathfrak a_s\Lambda_s\subset F_s,\quad
              \frac{\log N\mathfrak a_s}{d_L}\le X.
 \tag{HB8}
\]
Set \(j=r+k\), and define
\[
 \Phi=\sum_s\operatorname{ldeg}_s(F_s)
                   -\tfrac12\log H+\tfrac12f_0r\log(4\pi).
\]
Equations (HB3), (HB4) and (HB7) prove
\[
 \Phi\le f_0jX-f_0r\log R_E
            +\sum_s\widehat{\deg}_n M_s
            +\tfrac12f_0r\log(4\pi).
 \tag{HB9}\label{eq:sr-hb-parts}
\]
In fact, the torus indices contribute
\(f_0\log H/d_E=\tfrac12\log H\).
This is the exact cancellation with the term in \(\Phi\).
Bounding that index separately would lose this cancellation.

\subsection{An explicit small-block slope bound}

The following numerical lemma is the needed abstract deduction
from slope and conjugate-degree bounds.
It does not use an unproved geometric height theorem.

\begin{srlemma}\label{lem:sr-hb-small-blocks}
Let \(B\) be a Hermitian lattice of rank \(m\).
Suppose its generic space is the orthogonal sum of subspaces
\(U_s,W_s\), for \(f_0\) pairs.  Use their saturated integral
intersection lattices in the following degrees.
Let \(\Lambda_0\ge1\), \(j,n\ge1\), \(C\ge0\), and suppose
\[
 \begin{gathered}
 m\le f_0n\le e^{\Lambda_0},\qquad
 \operatorname{rk}U_s\le j,\\
 |\widehat{\deg}_nU_s-\widehat{\deg}_nW_s|
                                    \le Cn\Lambda_0,\\
 \mu_{\max}(B)\le C\log(2+m+\max(-\widehat{\deg}_nB,0)).
 \end{gathered}                                      \tag{HB10}
\]
Then, with \(C_a=C(2+\log(3+C))\),
\[
              \sum_s\widehat{\deg}_nU_s
                                \le C_a f_0j\Lambda_0.
 \tag{HB11}\label{eq:sr-hb-small-bound}
\]
\end{srlemma}

\begin{proof}
Write \(T=\sum_s\widehat{\deg}_nU_s\).
The direct sum of all the intersection lattices is a full
sublattice of \(B\); orthogonality makes its degree the sum
of their degrees.  Its index in \(B\) has nonnegative logarithm.
Therefore
\[
 \widehat{\deg}_nB\ge
      \sum_s(\widehat{\deg}_nU_s+\widehat{\deg}_nW_s)
                                  \ge2T-Cf_0n\Lambda_0.
\]
If \(T\le0\), the conclusion follows at once.
Otherwise \(h=\max(-\widehat{\deg}_nB,0)
\le Cf_0n\Lambda_0\le Ce^{\Lambda_0}\Lambda_0\).
Thus
\[
 \begin{split}
 \log(2+m+h)
 &\le\Lambda_0+\log(3+C\Lambda_0)\\
 &\le2\Lambda_0+\log(3+C)
 \le(2+\log(3+C))\Lambda_0.
 \end{split}
\]
The middle inequality uses
\(\Lambda_0\le e^{\Lambda_0}\) and \(1\le e^{\Lambda_0}\).
By definition of maximal slope, every nonzero saturated
\(U_s\) has degree at most its rank times \(\mu_{\max}(B)\).
The same statement is zero for rank zero.
Since \(T>0\), the maximal slope is positive.
Summing the rank bounds now gives (HB11).
\end{proof}

For the application, the spaces \(f(e_sV)\) have integral
image lattices \(M_s\) contained with finite index in the
corresponding saturated \(U_s\).
They have the same induced metrics, so
\(\widehat{\deg}_nM_s\le\widehat{\deg}_nU_s\).
The orthogonal \(E\)-action supplies the decomposition used
in (HB10).  A bound on its conjugate pair degrees and the
general Hodge maximal-slope bound are the two substantive
geometric inputs to that lemma; its conclusion follows
from exactly the displayed inequalities.

\subsection{The final numerical budget}

Suppose \(B_0\ge16\), \(R_E\ge1\), and put
\[
 \Lambda_0=\log B_0+\log R_E+|L'(1,\chi_d)/L(1,\chi_d)|.
\]
Suppose the prime set \(S\) and the actual labelled comparison
satisfy, with constants chosen before the host and point,
\[
 \sum_{p\in S}\frac{\log p}{p-1}\le C_1\log B_0,\qquad
 X=2\sum_{p\in S}\frac{\log p}{p-1}+2\log R_E,
 \tag{HB12}
\]
and suppose the small-block hypotheses hold with
\(f_0n\le B_0\) and the same \(\Lambda_0\).
Then \(X\le(2C_1+2)\Lambda_0\).
Moreover \(\log(4\pi)<\log16\le\Lambda_0\), and \(r\le j\).
Substitute (HB11) in (HB9), discard the nonpositive term
\(-f_0r\log R_E\), and obtain
\[
 \Phi\le(2C_1+2+C_a+\tfrac12)f_0j\Lambda_0
                     \le C_h f_0j\Lambda_0,\qquad
                 C_h=2C_1+C_a+3.
 \tag{HB13}\label{eq:sr-hb-final-budget}
\]
Every constant in (HB13) is chosen before \(E,L,B_0\), the
integral point and the individual label.

For clarity, (HB13) is the full determinant and numerical
assembly, not a claim that its geometric hypotheses have
already been discharged.  For the actual model one must
still establish (HB8), the genuine Hodge-metric identification,
the integral action on the abelian quotient, the conjugate
degree bound in (HB10), and the common compactification and
freeness conditions permitting all these statements to be
used at the same point.  Their source interfaces are,
respectively, \srcid{AbelianPartIntegralPoint},
\srcid{IsLinPtOf}, \srcid{exists_action},
\srcid{FlatSat}, and the admitted compact-datum intersection
conditions.  This chapter supplies the elementary conclusions
called \srcid{phiDeg_le_parts}, \srcid{abel_abstract}, and
\srcid{small_blocks_numbers} in the
\srcid{HeightComparisonGlobalModel} assembly.
It does not infer those model-dependent premises from their names.
